\scanpage{202}
par suite de type fini sur $k$. Soit en effet
$x$ un pt de $G_0$ rat sur $k$, il suffit de
prouver que $x\in U^{\unsure{2|?}}$ [N.B. $U^2$ est constructible
qu. cpct dans $G_0$] i.e. \deleted{$xU\cap U=$} l'existence
de $y$ et $z\in U(k)$ tels que $yz=x$, i.e.
$y=xz^{-1}$; ce qui signifie que $(U\cap xU)(k)\ne\varnothing$,
or $U$ et $xU$ sont des ouverts non vides
dans $G_0$, \deleted{\unsure{donc}}. Donc $G_0$ est de type
fini sur $k$. S'il n'était irréductible, \uline{tt}
point fermé de $G_0$ appartiendrait à \deleted{une
composante} deux composantes irréductibles
distinctes (par translation) absurde. On voit
aussi \deleted{\unsure{de suite}} 3° \deleted{\unsure{que}}\unsure{par|pour} \unsure{transp\textsuperscript{t}|transpos|transport} de
\deleted{Considérons la projection $G_k$}
structure [il suffit de considérer les
composantes connexes des pts fermés]
d'où aussi 1°. \deleted{\unsure{On obtient ainsi}} \inserted{\deleted{\unsure{Cela}} Prenant $U=G_0$, on}
voit que $G_0$ est un \uuline{sous-groupe}. Il est invariant
car \ldots

\deleted{Pour} Soit $k$ quelconque. \deleted{Soit} Il
suffit encore d'étudier les composantes
connexes \inserted{$U$} des pts \uline{fermés} $x$. \deleted{\unsure{Un tel}
pt devient rationnel sur une extension
finie $k'$, donc la composante connexe
\unsure{se décompose|décompose} en un nb fini}
% archive sheet number: 201
